\begin{lexlean}{Uorc.Page}
\title{Parameters}
\begin{semanticmodule}
structure{
  name: Workbench,
  type_parameters: [],
  parameters: [],
  fields: [
    {name: title, type: String},
    {name: heading, type: String},
    {name: purpose, type: String},
    {name: audience, type: String},
    {name: privacy, type: String},
    {name: guidance, type: (List String)},
    {name: initialOperation, type: String},
    {name: colorScheme, type: String},
    {name: density, type: String},
    {name: width, type: String}
  ]
}

structure{
  name: Labels,
  type_parameters: [],
  parameters: [],
  fields: [
    {name: skipToWorkbench, type: String},
    {name: workbench, type: String},
    {name: operations, type: String},
    {name: examples, type: String},
    {name: examplesHelp, type: String},
    {name: usage, type: String},
    {name: links, type: String},
    {name: noscript, type: String},
    {name: version, type: String},
    {name: progress, type: String},
    {name: download, type: String},
    {name: restored, type: String},
    {name: cancelled, type: String},
    {name: memory, type: String}
  ]
}

structure{
  name: Stakeholder,
  type_parameters: [],
  parameters: [],
  fields: [{name: name, type: String}, {name: text, type: String}]
}

structure{
  name: Feature,
  type_parameters: [],
  parameters: [],
  fields: [
    {name: statement, type: String},
    {name: calls, type: (List String)},
    {name: acceptance, type: (List String)}
  ]
}

structure{
  name: Journey,
  type_parameters: [],
  parameters: [],
  fields: [{name: actor, type: String}, {name: steps, type: (List String)}]
}

structure{
  name: UmlState,
  type_parameters: [],
  parameters: [],
  fields: [{name: state, type: String}, {name: initial, type: Bool}, {name: status, type: String}]
}

structure{
  name: UmlTransition,
  type_parameters: [],
  parameters: [],
  fields: [
    {name: order, type: Nat},
    {name: source, type: String},
    {name: target, type: String},
    {name: trigger, type: String}
  ]
}

structure{
  name: Operation,
  type_parameters: [],
  parameters: [],
  fields: [
    {name: id, type: String},
    {name: order, type: Nat},
    {name: label, type: String},
    {name: summary, type: String},
    {name: help, type: String},
    {name: fields, type: (List String)},
    {name: runLabel, type: String},
    {name: running, type: String},
    {name: result, type: String}
  ]
}

structure{
  name: Field,
  type_parameters: [],
  parameters: [],
  fields: [
    {name: id, type: String},
    {name: label, type: String},
    {name: kind, type: String},
    {name: help, type: String},
    {name: required, type: Bool},
    {name: replaces, type: String}
  ]
}

structure{
  name: Control,
  type_parameters: [],
  parameters: [],
  fields: [
    {name: id, type: String},
    {name: label, type: String},
    {name: trigger, type: String},
    {name: enabledIn, type: (List String)}
  ]
}

structure{
  name: Example,
  type_parameters: [],
  parameters: [],
  fields: [
    {name: id, type: String},
    {name: order, type: Nat},
    {name: label, type: String},
    {name: description, type: String},
    {name: input, type: String},
    {name: bytes, type: Nat},
    {name: expectedResponse, type: String}
  ]
}

structure{
  name: Column,
  type_parameters: [],
  parameters: [],
  fields: [
    {name: table, type: String},
    {name: order, type: Nat},
    {name: key, type: String},
    {name: label, type: String}
  ]
}

structure{
  name: ErrorHint,
  type_parameters: [],
  parameters: [],
  fields: [
    {name: code, type: String},
    {name: title, type: String},
    {name: explanation, type: String},
    {name: recovery, type: String},
    {name: fields, type: (List String)}
  ]
}

structure{
  name: Section,
  type_parameters: [],
  parameters: [],
  fields: [
    {name: id, type: String},
    {name: order, type: Nat},
    {name: title, type: String},
    {name: paragraphs, type: (List String)}
  ]
}

structure{
  name: Usage,
  type_parameters: [],
  parameters: [],
  fields: [
    {name: order, type: Nat},
    {name: title, type: String},
    {name: language, type: String},
    {name: code, type: (List String)}
  ]
}

structure{
  name: Link,
  type_parameters: [],
  parameters: [],
  fields: [{name: order, type: Nat}, {name: label, type: String}, {name: href, type: String}]
}

structure{
  name: QualityRequirement,
  type_parameters: [],
  parameters: [],
  fields: [
    {name: characteristic, type: String},
    {name: subcharacteristic, type: String},
    {name: text, type: String}
  ]
}

structure{
  name: QualityMeasure,
  type_parameters: [],
  parameters: [],
  fields: [
    {name: requirement, type: String},
    {name: quantity, type: String},
    {name: unit, type: String},
    {name: threshold, type: Nat},
    {name: disposition, type: String}
  ]
}

structure{
  name: DecodeHost,
  type_parameters: [],
  parameters: [],
  fields: [
    {name: workersPerGiB, type: Nat},
    {name: coreLimit, type: Nat},
    {name: unreportedGiB, type: Nat},
    {name: keepsTables, type: Bool},
    {name: text, type: String}
  ]
}

definition{
  name: workbench,
  parameters: [],
  result: Workbench,
  body: (record Workbench {title := "uorc client", heading := "uorc", purpose := "Compress files and folders into .uorc archives, and restore, inspect or verify any .uorc archive, in this browser.", audience := "Anyone with files to compress or an archive to open, and developers judging the uorc format before installing the command line.", privacy := "Your files stay on this device: the page reads them, runs the uorc WebAssembly client in a background worker and hands the result back as a download. Nothing is uploaded.", guidance := (cons "Pick what to do: Compress, Decompress, Inspect or Verify." (cons "Choose the files, folder or archive it asks for." (cons "Press the button under the fields, then save the result. Cancel stops a running operation at once." (nil String)))), initialOperation := "compress", colorScheme := "system", density := "comfortable", width := "balanced"})
}

definition{
  name: labels,
  parameters: [],
  result: Labels,
  body: (record Labels {skipToWorkbench := "Skip to the workbench", workbench := "Use uorc here", operations := "Operation", examples := "Examples", examplesHelp := "Each example compresses a built-in input in this page, decompresses the archive again and checks that the bytes are identical.", usage := "Use uorc elsewhere", links := "More", noscript := "This page runs the uorc client with JavaScript and WebAssembly; enable JavaScript to use it.", version := "Release", progress := "Progress", download := "Save", restored := "Restored files", cancelled := "Cancelled. Ready.", memory := "WebAssembly memory"})
}

definition{
  name: visitor,
  parameters: [],
  result: Stakeholder,
  body: (record Stakeholder {name := "Visitor", text := "Someone who wants to compress files or open a .uorc archive without installing anything."})
}

definition{
  name: evaluator,
  parameters: [],
  result: Stakeholder,
  body: (record Stakeholder {name := "Evaluator", text := "A developer judging the uorc format and client: what an archive holds, whether it decodes, how the page and the command line agree."})
}

definition{
  name: compressFile,
  parameters: [],
  result: Feature,
  body: (record Feature {statement := "Compress one file of any size into a .uorc archive with exactly the bytes the uorc command line writes for it.", calls := (cons "uorc.compressBlob" (nil String)), acceptance := (cons "compress-file" (cons "compress-matches-library" (nil String)))})
}

definition{
  name: compressFolder,
  parameters: [],
  result: Feature,
  body: (record Feature {statement := "Compress several files, or one folder with its relative paths, into one archive of ordered objects.", calls := (cons "uorc.compressTreeBlob" (nil String)), acceptance := (cons "compress-files" (cons "compress-folder" (nil String)))})
}

definition{
  name: decompressArchive,
  parameters: [],
  result: Feature,
  body: (record Feature {statement := "Restore any .uorc archive this page or the command line wrote, including archives of earlier releases and streams of several frames with reference frames, and offer every restored file as a download.", calls := (cons "uorc.decompressBlob" (cons "uorc.decompressTreeBlob" (nil String))), acceptance := (cons "decompress-file" (cons "decompress-reference-frames" (cons "decompress-folder" (cons "decompress-released" (nil String)))))})
}

definition{
  name: inspectArchive,
  parameters: [],
  result: Feature,
  body: (record Feature {statement := "Read every archive header and integrity seal of a stream without decoding it, as a table.", calls := (cons "uorc.inspect" (nil String)), acceptance := (cons "inspect-archive" (nil String))})
}

definition{
  name: verifyArchive,
  parameters: [],
  result: Feature,
  body: (record Feature {statement := "Decode an archive without keeping its output and, given the original file, check that the bytes are identical.", calls := (cons "uorc.verify" (nil String)), acceptance := (cons "verify-archive" (cons "verify-original" (cons "verify-original-differs" (nil String))))})
}

definition{
  name: cancelOperation,
  parameters: [],
  result: Feature,
  body: (record Feature {statement := "Stop a running operation at once, discard its partial result and keep the page ready for the next one.", calls := (cons "Worker.terminate" (nil String)), acceptance := (cons "cancel-running" (nil String))})
}

definition{
  name: runExamples,
  parameters: [],
  result: Feature,
  body: (record Feature {statement := "Run built-in inputs through compress and decompress and show their sizes, times and whether the round trip is identical.", calls := (cons "uorc.compress" (cons "uorc.decompress" (nil String))), acceptance := (cons "examples-round-trip" (nil String))})
}

definition{
  name: explainFailures,
  parameters: [],
  result: Feature,
  body: (record Feature {statement := "Explain every failure with its cause and the next step, and keep the page usable afterwards.", calls := (cons "uorc.ERRORS" (nil String)), acceptance := (cons "error-decode" (cons "error-mismatch" (cons "error-stopped" (cons "error-load" (nil String)))))})
}

definition{
  name: compressThenRestore,
  parameters: [],
  result: Journey,
  body: (record Journey {actor := "Uorc.Page::visitor", steps := (cons "Uorc.Page::compressFile" (cons "Uorc.Page::decompressArchive" (cons "Uorc.Page::cancelOperation" (nil String))))})
}

definition{
  name: packFolder,
  parameters: [],
  result: Journey,
  body: (record Journey {actor := "Uorc.Page::visitor", steps := (cons "Uorc.Page::compressFolder" (cons "Uorc.Page::decompressArchive" (nil String)))})
}

definition{
  name: judgeFormat,
  parameters: [],
  result: Journey,
  body: (record Journey {actor := "Uorc.Page::evaluator", steps := (cons "Uorc.Page::runExamples" (cons "Uorc.Page::inspectArchive" (cons "Uorc.Page::verifyArchive" (cons "Uorc.Page::explainFailures" (nil String)))))})
}

definition{
  name: stateLoading,
  parameters: [],
  result: UmlState,
  body: (record UmlState {state := "loading", initial := true, status := "Loading the uorc WebAssembly client…"})
}

definition{
  name: stateUnavailable,
  parameters: [],
  result: UmlState,
  body: (record UmlState {state := "unavailable", initial := false, status := "The uorc client could not start in this browser."})
}

definition{
  name: stateReady,
  parameters: [],
  result: UmlState,
  body: (record UmlState {state := "ready", initial := false, status := "Ready."})
}

definition{
  name: stateRunning,
  parameters: [],
  result: UmlState,
  body: (record UmlState {state := "running", initial := false, status := "Working…"})
}

definition{
  name: stateCancelling,
  parameters: [],
  result: UmlState,
  body: (record UmlState {state := "cancelling", initial := false, status := "Cancelling…"})
}

definition{
  name: stateDone,
  parameters: [],
  result: UmlState,
  body: (record UmlState {state := "done", initial := false, status := "Done."})
}

definition{
  name: stateFailed,
  parameters: [],
  result: UmlState,
  body: (record UmlState {state := "failed", initial := false, status := "Failed."})
}

definition{
  name: whenLoaded,
  parameters: [],
  result: UmlTransition,
  body: (record UmlTransition {order := 1, source := "loading", target := "ready", trigger := "loaded"})
}

definition{
  name: whenLoadFails,
  parameters: [],
  result: UmlTransition,
  body: (record UmlTransition {order := 2, source := "loading", target := "unavailable", trigger := "loadFailed"})
}

definition{
  name: runFromReady,
  parameters: [],
  result: UmlTransition,
  body: (record UmlTransition {order := 3, source := "ready", target := "running", trigger := "run"})
}

definition{
  name: whenFinished,
  parameters: [],
  result: UmlTransition,
  body: (record UmlTransition {order := 4, source := "running", target := "done", trigger := "finished"})
}

definition{
  name: whenFailed,
  parameters: [],
  result: UmlTransition,
  body: (record UmlTransition {order := 5, source := "running", target := "failed", trigger := "failed"})
}

definition{
  name: cancelRunning,
  parameters: [],
  result: UmlTransition,
  body: (record UmlTransition {order := 6, source := "running", target := "cancelling", trigger := "cancel"})
}

definition{
  name: whenRestarted,
  parameters: [],
  result: UmlTransition,
  body: (record UmlTransition {order := 7, source := "cancelling", target := "ready", trigger := "restarted"})
}

definition{
  name: runFromDone,
  parameters: [],
  result: UmlTransition,
  body: (record UmlTransition {order := 8, source := "done", target := "running", trigger := "run"})
}

definition{
  name: runFromFailed,
  parameters: [],
  result: UmlTransition,
  body: (record UmlTransition {order := 9, source := "failed", target := "running", trigger := "run"})
}

definition{
  name: editAfterDone,
  parameters: [],
  result: UmlTransition,
  body: (record UmlTransition {order := 10, source := "done", target := "ready", trigger := "edit"})
}

definition{
  name: editAfterFailure,
  parameters: [],
  result: UmlTransition,
  body: (record UmlTransition {order := 11, source := "failed", target := "ready", trigger := "edit"})
}

definition{
  name: operationCompress,
  parameters: [],
  result: Operation,
  body: (record Operation {id := "compress", order := 1, label := "Compress", summary := "Make a .uorc archive from files or a folder.", help := "One file becomes NAME.uorc with exactly the bytes the uorc command line writes. Several files or a folder become one archive that keeps their relative paths.", fields := (cons "files" (cons "folder" (nil String))), runLabel := "Compress", running := "Compressing", result := "The archive, ready to save, with its size against the input."})
}

definition{
  name: operationDecompress,
  parameters: [],
  result: Operation,
  body: (record Operation {id := "decompress", order := 2, label := "Decompress", summary := "Restore the files a .uorc archive holds.", help := "Any archive the page or the command line wrote. An archive of one file restores that file; an archive of several files or a folder lists every file to save.", fields := (cons "archive" (nil String)), runLabel := "Decompress", running := "Decompressing", result := "Every restored file, ready to save."})
}

definition{
  name: operationInspect,
  parameters: [],
  result: Operation,
  body: (record Operation {id := "inspect", order := 3, label := "Inspect", summary := "See what an archive holds without decoding it.", help := "Reads each archive header of the stream and checks its integrity seal. Fast for any size.", fields := (cons "archive" (nil String)), runLabel := "Inspect", running := "Inspecting", result := "One row per archive of the stream."})
}

definition{
  name: operationVerify,
  parameters: [],
  result: Operation,
  body: (record Operation {id := "verify", order := 4, label := "Verify", summary := "Check that an archive decodes, and optionally that it reproduces an original.", help := "Decodes the whole archive without saving anything. Give the original file too and every byte is compared.", fields := (cons "archive" (cons "original" (nil String))), runLabel := "Verify", running := "Verifying", result := "Whether the archive decodes and, with an original, whether the bytes are identical."})
}

definition{
  name: fieldFiles,
  parameters: [],
  result: Field,
  body: (record Field {id := "files", label := "Files", kind := "files", help := "One or more files. Choosing files clears a chosen folder.", required := true, replaces := "folder"})
}

definition{
  name: fieldFolder,
  parameters: [],
  result: Field,
  body: (record Field {id := "folder", label := "Or a folder", kind := "folder", help := "Every file under one folder, with its relative paths. Choosing a folder clears chosen files.", required := true, replaces := "files"})
}

definition{
  name: fieldArchive,
  parameters: [],
  result: Field,
  body: (record Field {id := "archive", label := "Archive", kind := "archive", help := "A .uorc file written by this page or the uorc command line.", required := true, replaces := ""})
}

definition{
  name: fieldOriginal,
  parameters: [],
  result: Field,
  body: (record Field {id := "original", label := "Original (optional)", kind := "file", help := "The file the archive should reproduce. Leave empty to check only that the archive decodes.", required := false, replaces := ""})
}

definition{
  name: controlRun,
  parameters: [],
  result: Control,
  body: (record Control {id := "run", label := "Run the chosen operation", trigger := "run", enabledIn := (cons "ready" (cons "done" (cons "failed" (nil String))))})
}

definition{
  name: controlCancel,
  parameters: [],
  result: Control,
  body: (record Control {id := "cancel", label := "Cancel", trigger := "cancel", enabledIn := (cons "running" (nil String))})
}

definition{
  name: controlOperation,
  parameters: [],
  result: Control,
  body: (record Control {id := "operation", label := "Operation", trigger := "edit", enabledIn := (cons "ready" (cons "done" (cons "failed" (nil String))))})
}

definition{
  name: controlField,
  parameters: [],
  result: Control,
  body: (record Control {id := "field", label := "Field", trigger := "edit", enabledIn := (cons "ready" (cons "done" (cons "failed" (nil String))))})
}

definition{
  name: controlExample,
  parameters: [],
  result: Control,
  body: (record Control {id := "example", label := "Example", trigger := "run", enabledIn := (cons "ready" (cons "done" (cons "failed" (nil String))))})
}

definition{
  name: exampleText,
  parameters: [],
  result: Example,
  body: (record Example {id := "text", order := 1, label := "English text", description := "This page's own guidance and purpose, repeated: prose compresses well.", input := "text", bytes := 16384, expectedResponse := "identical"})
}

definition{
  name: exampleZeros,
  parameters: [],
  result: Example,
  body: (record Example {id := "zeros", order := 2, label := "Zero bytes", description := "One run of a single byte value.", input := "zeros", bytes := 65536, expectedResponse := "identical"})
}

definition{
  name: exampleCounter,
  parameters: [],
  result: Example,
  body: (record Example {id := "counter", order := 3, label := "Counter bytes", description := "Every byte value 0 to 255 in order, repeated.", input := "counter", bytes := 4096, expectedResponse := "identical"})
}

definition{
  name: exampleRandom,
  parameters: [],
  result: Example,
  body: (record Example {id := "random", order := 4, label := "Random bytes", description := "Bytes from the browser's random source: they do not compress, and the archive stores them as they are.", input := "random", bytes := 16384, expectedResponse := "identical"})
}

definition{
  name: exampleEmpty,
  parameters: [],
  result: Example,
  body: (record Example {id := "empty", order := 5, label := "Empty input", description := "No bytes at all.", input := "empty", bytes := 0, expectedResponse := "identical"})
}

definition{
  name: columnArchive,
  parameters: [],
  result: Column,
  body: (record Column {table := "inspect", order := 1, key := "index", label := "Archive"})
}

definition{
  name: columnRepresentation,
  parameters: [],
  result: Column,
  body: (record Column {table := "inspect", order := 2, key := "representation", label := "Coding"})
}

definition{
  name: columnObjects,
  parameters: [],
  result: Column,
  body: (record Column {table := "inspect", order := 3, key := "objects", label := "Objects"})
}

definition{
  name: columnOutput,
  parameters: [],
  result: Column,
  body: (record Column {table := "inspect", order := 4, key := "output_bytes", label := "Restored bytes"})
}

definition{
  name: columnArchiveBytes,
  parameters: [],
  result: Column,
  body: (record Column {table := "inspect", order := 5, key := "archive_bytes", label := "Archive bytes"})
}

definition{
  name: columnIntegrity,
  parameters: [],
  result: Column,
  body: (record Column {table := "inspect", order := 6, key := "integrity_ok", label := "Integrity seal"})
}

definition{
  name: columnExampleInput,
  parameters: [],
  result: Column,
  body: (record Column {table := "examples", order := 1, key := "label", label := "Input"})
}

definition{
  name: columnExampleBytes,
  parameters: [],
  result: Column,
  body: (record Column {table := "examples", order := 2, key := "bytes", label := "Bytes"})
}

definition{
  name: columnExampleArchive,
  parameters: [],
  result: Column,
  body: (record Column {table := "examples", order := 3, key := "archive", label := "Archive bytes"})
}

definition{
  name: columnExampleRoundTrip,
  parameters: [],
  result: Column,
  body: (record Column {table := "examples", order := 4, key := "roundTrip", label := "Round trip"})
}

definition{
  name: columnExampleTime,
  parameters: [],
  result: Column,
  body: (record Column {table := "examples", order := 5, key := "seconds", label := "Seconds"})
}

definition{
  name: errorBudget,
  parameters: [],
  result: ErrorHint,
  body: (record ErrorHint {code := "budget", title := "The archive needs more decoding work than allowed", explanation := "Decoding stops at a fixed work limit so a damaged or hostile archive cannot run forever. Archives uorc writes stay far below it.", recovery := "Check that the file is a .uorc archive from a trusted source, or decode it with the uorc command line and a larger --work value.", fields := (cons "archive" (nil String))})
}

definition{
  name: errorOverflow,
  parameters: [],
  result: ErrorHint,
  body: (record ErrorHint {code := "overflow", title := "The archive asks for numbers too large to decode", explanation := "A length or count in the archive does not fit in 64 bits, which no archive uorc writes does.", recovery := "The file is damaged or is not a uorc archive. Use the original archive.", fields := (cons "archive" (nil String))})
}

definition{
  name: errorDecode,
  parameters: [],
  result: ErrorHint,
  body: (record ErrorHint {code := "decode", title := "This is not a uorc archive, or it is damaged", explanation := "The file does not decode as a stream of .uorc archives: its header, body or integrity seal does not check.", recovery := "Choose the .uorc file itself. To make an archive from this file, use Compress.", fields := (cons "archive" (nil String))})
}

definition{
  name: errorMismatch,
  parameters: [],
  result: ErrorHint,
  body: (record ErrorHint {code := "mismatch", title := "The archive does not reproduce the original", explanation := "The archive decodes, but its bytes differ from the original you chose.", recovery := "Check that you chose the original the archive was made from.", fields := (cons "original" (nil String))})
}

definition{
  name: errorStopped,
  parameters: [],
  result: ErrorHint,
  body: (record ErrorHint {code := "stopped", title := "The browser stopped the operation", explanation := "The background worker running uorc ended before finishing, most often because the browser refused it more memory.", recovery := "Try again; for very large inputs use the uorc command line, which streams with bounded memory.", fields := (nil String)})
}

definition{
  name: errorLoad,
  parameters: [],
  result: ErrorHint,
  body: (record ErrorHint {code := "load", title := "The uorc client could not start", explanation := "This browser could not load the uorc WebAssembly module in a worker. It needs module workers and WebAssembly SIMD (current Chrome, Edge, Firefox and Safari 16.4 or later).", recovery := "Reload the page, or open it in a current browser.", fields := (nil String)})
}

definition{
  name: sectionAbout,
  parameters: [],
  result: Section,
  body: (record Section {id := "about", order := 1, title := "What the page does", paragraphs := (cons "uorc is Universal Object Reference Compression. Each archive is an exact executable object reference: it carries the program that rebuilds the original, and an integrity seal over itself." (cons "This page runs the same uorc client as the command line, compiled to WebAssembly. Compress writes the same bytes, Decompress reads every archive the command line writes, and Verify checks an archive without saving anything." (cons "Large inputs stream through 8 MiB frames, so a file of any size goes through a few frames of memory at a time. Several files or a folder are compressed together in memory as one archive of ordered objects: a manifest of relative paths, then each file." (nil String))))})
}

definition{
  name: usageCommandLine,
  parameters: [],
  result: Usage,
  body: (record Usage {order := 1, title := "Command line", language := "shell", code := (cons "uorc compress --stats notes.txt      # writes notes.txt.uorc" (cons "uorc decompress notes.txt.uorc -o notes.restored.txt" (cons "cat data.bin | uorc compress - > data.uorc   # any size: streams through 8 MiB frames" (cons "uorc compress --threads 4 big.tar     # frames compress in parallel" (cons "uorc decompress --threads 4 big.tar.uorc   # blocks of frames decode in parallel" (cons "uorc compress photos/                 # writes photos.uorc (a whole tree)" (cons "uorc decompress photos.uorc -o restore/   # --resume finishes an interrupted restore" (cons "uorc inspect notes.txt.uorc           # header and integrity, no decoding" (cons "uorc verify notes.txt.uorc notes.txt  # decode, check, compare" (cons "# Outputs are written to a temporary file and renamed into place; --force replaces." (cons "# Exit status: 0 ok, 2 usage, 3 I/O, 4 invalid archive or integrity, 5 step budget," (cons "# 6 overflow, 7 differs from the original." (nil String)))))))))))))})
}

definition{
  name: usageRust,
  parameters: [],
  result: Usage,
  body: (record Usage {order := 2, title := "Rust library", language := "rust", code := (cons "let archive = uorc::compress(&bytes, uorc::DEFAULT_STEPS)?;" (cons "let original = uorc::decompress(&archive, uorc::DEFAULT_DECODE_WORK)?;" (cons "assert_eq!(original, bytes);" (cons "let info = uorc::inspect(&archive)?;           // header + integrity, no decode" (cons "uorc::verify(&archive, Some(&bytes), uorc::DEFAULT_DECODE_WORK)?;" (cons "let many = uorc::compress_objects(&[a, b, c], uorc::DEFAULT_STEPS)?;" (cons "let parts = uorc::decompress_objects(&many, uorc::DEFAULT_DECODE_WORK)?;" (cons "// Or stream between any reader and writer, with bounded memory:" (cons "uorc::compress_stream(reader, writer, uorc::DEFAULT_STEPS, uorc::default_threads())?;" (nil String))))))))))})
}

definition{
  name: usageJavaScript,
  parameters: [],
  result: Usage,
  body: (record Usage {order := 3, title := "JavaScript (this page's module)", language := "javascript", code := (cons "import { load } from \"./uorc.js\";" (cons "const uorc = await load(\"./uorc_bg.wasm\");" (cons "const archive = uorc.compress(new TextEncoder().encode(\"hello hello hello\"));" (cons "const original = uorc.decompress(archive);" (cons "uorc.inspect(archive);                 // [{ representation, object_lengths, integrity_ok, ... }]" (cons "uorc.verify(archive, original);        // throws if it does not reproduce original" (cons "const bundle = uorc.compressFiles([{ path: \"a.txt\", data }, { path: \"d\", kind: \"directory\" }]);" (cons "const files = uorc.decompressFiles(bundle); // [{ path, kind, data }]" (cons "// Files of any size, one frame in memory at a time:" (cons "const blob = await uorc.compressBlob(file);" (cons "const back = await uorc.decompressBlob(blob);" (cons "// Folders of any size, the same way:" (cons "const tree = await uorc.compressTreeBlob([{ path: \"d\", kind: \"directory\" }, { path: \"d/a.txt\", blob: file }]);" (cons "const entries = await uorc.decompressTreeBlob(tree); // [{ path, kind, blob }]" (nil String)))))))))))))))})
}

definition{
  name: linkReleases,
  parameters: [],
  result: Link,
  body: (record Link {order := 1, label := "Releases: binaries, the library source and this WebAssembly module", href := "https://github.com/UOR-Foundation/compression/releases"})
}

definition{
  name: linkBenchmarks,
  parameters: [],
  result: Link,
  body: (record Link {order := 2, label := "Benchmarks for this release", href := "bench.html"})
}

definition{
  name: linkSource,
  parameters: [],
  result: Link,
  body: (record Link {order := 3, label := "Source code", href := "https://github.com/UOR-Foundation/compression"})
}

definition{
  name: responsiveWhileRunning,
  parameters: [],
  result: QualityRequirement,
  body: (record QualityRequirement {characteristic := "performance efficiency", subcharacteristic := "time behaviour", text := "The page answers input while an operation runs: the codec runs in a worker, never on the page's thread."})
}

definition{
  name: cancelAtOnce,
  parameters: [],
  result: QualityRequirement,
  body: (record QualityRequirement {characteristic := "interaction capability", subcharacteristic := "user error protection", text := "Cancel stops a running operation at once, whatever the input size."})
}

definition{
  name: startQuickly,
  parameters: [],
  result: QualityRequirement,
  body: (record QualityRequirement {characteristic := "performance efficiency", subcharacteristic := "time behaviour", text := "The page is ready to use soon after it loads."})
}

definition{
  name: sameBytesAsLibrary,
  parameters: [],
  result: QualityRequirement,
  body: (record QualityRequirement {characteristic := "functional suitability", subcharacteristic := "functional correctness", text := "An archive the page writes is byte-identical to the archive the uorc library and command line write for the same input."})
}

definition{
  name: boundedMemory,
  parameters: [],
  result: QualityRequirement,
  body: (record QualityRequirement {characteristic := "performance efficiency", subcharacteristic := "resource utilization", text := "Compressing or decompressing one file holds a few frames in memory, not the whole file, and one worker holds at most its share of the device (Uorc.Page::pageWorker). During a frame the worker holds one set of the coder's tables, the window's frames read in place and the output; between frames it keeps the tables as spares for the next frame, so the next frame allocates nothing. The encoder counts all of these against the worker's share before it grows a frame's tables. At v0.4.0, decoding distinct 8 MiB frames peaked at 226 MiB for three frames and 250 MiB for six or twelve; what the worker retains beyond the decoder is follow-up for the page."})
}

definition{
  name: statusMessages,
  parameters: [],
  result: QualityRequirement,
  body: (record QualityRequirement {characteristic := "interaction capability", subcharacteristic := "inclusivity", text := "WCAG 2.2 4.1.3 and 4.1.2: progress and results are announced through one status region, and every control has an accessible name."})
}

definition{
  name: browserThroughput,
  parameters: [],
  result: QualityRequirement,
  body: (record QualityRequirement {characteristic := "performance efficiency", subcharacteristic := "time behaviour", text := "Compress and decompress throughput in the browser on 10 MiB of text, measured and reported. Measured on the same 10 MiB (4 vCPU, headless Chromium): one worker compressed it at 168 KB/s and decompressed it at 563 KB/s; with the worker pool the page compresses it at 369 KB/s and decompresses it at 561 KB/s. Decoding this input stays on one worker, as it does natively, because each of its frames is a window frame whose model reads the frames before it."})
}

definition{
  name: useTheCores,
  parameters: [],
  result: QualityRequirement,
  body: (record QualityRequirement {characteristic := "performance efficiency", subcharacteristic := "resource utilization", text := "Compressing uses the visitor's cores: a pool of workers codes each frame's two scheduled strategies side by side, and frames side by side, and the archive is still the library's."})
}

definition{
  name: progressWithinFrame,
  parameters: [],
  result: QualityRequirement,
  body: (record QualityRequirement {characteristic := "interaction capability", subcharacteristic := "self-descriptiveness", text := "Progress moves while a frame is coded, not only between frames, so a long operation never looks stalled."})
}

definition{
  name: coresUnused,
  parameters: [],
  result: QualityMeasure,
  body: (record QualityMeasure {requirement := "Uorc.Page::useTheCores", quantity := "cores the page may use (at most 16, and two per GiB of device memory) without a pool worker", unit := "cores", threshold := 0, disposition := "gate"})
}

definition{
  name: progressGap,
  parameters: [],
  result: QualityMeasure,
  body: (record QualityMeasure {requirement := "Uorc.Page::progressWithinFrame", quantity := "longest time the progress bar stands still while an operation runs", unit := "ms", threshold := 2000, disposition := "gate"})
}

definition{
  name: longestPageTask,
  parameters: [],
  result: QualityMeasure,
  body: (record QualityMeasure {requirement := "Uorc.Page::responsiveWhileRunning", quantity := "longest task on the page's thread while an operation runs", unit := "ms", threshold := 200, disposition := "gate"})
}

definition{
  name: cancelLatency,
  parameters: [],
  result: QualityMeasure,
  body: (record QualityMeasure {requirement := "Uorc.Page::cancelAtOnce", quantity := "time from Cancel to the page being ready again", unit := "ms", threshold := 3000, disposition := "gate"})
}

definition{
  name: readyTime,
  parameters: [],
  result: QualityMeasure,
  body: (record QualityMeasure {requirement := "Uorc.Page::startQuickly", quantity := "time from navigation to the page being ready", unit := "ms", threshold := 5000, disposition := "gate"})
}

definition{
  name: differingArchives,
  parameters: [],
  result: QualityMeasure,
  body: (record QualityMeasure {requirement := "Uorc.Page::sameBytesAsLibrary", quantity := "page archives that differ from the library's for the same input", unit := "archives", threshold := 0, disposition := "gate"})
}

definition{
  name: pageWorker,
  parameters: [],
  result: DecodeHost,
  body: (record DecodeHost {workersPerGiB := 2, coreLimit := 16, unreportedGiB := 2, keepsTables := true, text := "The page's host rule, as a function of the device the browser reports: it starts one worker per core, at most coreLimit and at most workersPerGiB per GiB of navigator.deviceMemory (unreportedGiB when the browser reports none), so each worker may hold 1 GiB / workersPerGiB of WebAssembly memory (512 MiB). That is the decode host every archive is written for by default (uorc performance::PAGE_WORKER, Client.hostMemory): a frame's tables grow only while its decoder, the spare set of tables the worker keeps between frames, and the frame's window and output fit it. One worker decodes a frame. workersPerGiB, coreLimit and unreportedGiB are stated, not yet derived from the device."})
}

definition{
  name: workerMemory,
  parameters: [],
  result: QualityMeasure,
  body: (record QualityMeasure {requirement := "Uorc.Page::boundedMemory", quantity := "WebAssembly memory after decompressing the 10 MiB text archive, at most one page worker's share (Uorc.Page::pageWorker: 1 GiB / workersPerGiB)", unit := "MiB", threshold := 512, disposition := "gate"})
}

definition{
  name: unnamedControls,
  parameters: [],
  result: QualityMeasure,
  body: (record QualityMeasure {requirement := "Uorc.Page::statusMessages", quantity := "controls without an accessible name", unit := "controls", threshold := 0, disposition := "gate"})
}

definition{
  name: compressThroughput,
  parameters: [],
  result: QualityMeasure,
  body: (record QualityMeasure {requirement := "Uorc.Page::browserThroughput", quantity := "compress throughput on 10 MiB of text, on every worker the page starts", unit := "KB/s", threshold := 0, disposition := "measured"})
}

definition{
  name: decompressThroughput,
  parameters: [],
  result: QualityMeasure,
  body: (record QualityMeasure {requirement := "Uorc.Page::browserThroughput", quantity := "decompress throughput on 10 MiB of text", unit := "KB/s", threshold := 0, disposition := "measured"})
}
\end{semanticmodule}
\end{lexlean}
